% letter-two-page.tex — layout stress: enough prose to span two pages, so a page
% break falls inside the body. Verifies that body paragraphs and the closing
% break across pages without a page-number folio appearing and without overfull
% boxes.
\documentclass{careerdossier-letter}
\CDossierSetup{
  name     = {Ada Lovelace},
  headline = {Data Scientist},
  email    = {ada.lovelace@example.com},
  phone    = {+1 416 555 0142},
  location = {Toronto, Ontario, Canada},
  linkedin = {linkedin.com/in/ada-lovelace},
  github   = {github.com/ada-lovelace},
}
\CDossierLetterSetup{
  date                   = {14 July 2026},
  recipient-name         = {Hiring Committee},
  recipient-title        = {Department of Applied Analytics},
  recipient-organization = {Northbridge Research Institute},
  recipient-address      = {123 Discovery Avenue\\Vancouver, British Columbia V6T 1Z4},
  subject                = {Application for the Senior Data Scientist position},
  salutation             = {Dear Members of the Hiring Committee,},
}
\begin{document}
\MakeCDossierLetterhead

I am writing to apply for the Senior Data Scientist position in computational
analytics at the Northbridge Research Institute. I recently completed a decade
of work spanning public-sector analytics, reporting platforms, and reproducible
research pipelines, and I am eager to contribute that experience to your group's
mission of building reliable, accessible scientific software.

My most recent work focused on developing scalable, reproducible reporting
services for high-volume public programs. I designed governed data pipelines,
implemented performance-critical processing routines, and collaborated closely
with program teams to translate operational requirements into maintainable
systems with measurable service standards. This work reduced duplicated
maintenance across teams, shortened incident investigation, and established
automated validation for data quality, reproducibility, and archival output.

Earlier in my career I built data-import services and validation tools for
public reporting datasets in a range of interchange formats. I improved
processing performance through batching, indexing, and query optimization while
preserving auditability, and I created internal libraries for schema validation,
error reporting, and automated documentation generation. Across these roles I
have consistently favoured tested, documented, and reproducible tools over
one-off analyses, because they are what allow a team to move quickly without
sacrificing reliability.

The position is especially appealing because of its emphasis on interdisciplinary
collaboration and dependable scientific software. My experience translating
analytical models into tested computational tools would allow me to contribute
quickly to your current projects, and I would welcome the opportunity to mentor
junior researchers and to support the group's documentation and reproducibility
practices. I am particularly interested in the institute's work on accessible
public correspondence platforms, where careful attention to content models,
required-content checks, reading order, and archival conformance mirrors the
standards I have applied throughout my own projects.

Beyond the technical work, I care deeply about making analytical systems
understandable to the people who depend on them. I have written operational
documentation, run internal training sessions, and defined reusable standards for
testing, release management, and technical writing. I believe that reliable
software and clear communication are two halves of the same responsibility, and I
try to hold both to the same standard in everything I deliver.

Thank you for considering my application. I have enclosed my curriculum vitae and
would be pleased to discuss how my background and interests align with the goals
of your group. I am available at your convenience for a conversation and can
provide references, code samples, and further documentation on request.

\MakeCDossierClosing
\end{document}
